% Wave-10 S1/20: the 2-divisibility orbit framework for N in {4,6,8}.
% Successor to wave8_l5_slides_boundary_N_5_7_8.tex.
% Primary sources: Mukai 1988 Invent. Math. 94, 183-221; Clery-Gritsenko
% 2008 (arXiv:0812.3962), 2013 (IJM 24:1350059); Cheng-Harvey-Paquette
% 2014 (arXiv:1406.5502). Slides: /tmp/automorphic-slides.txt.

\providecommand{\kBKM}{\kappa_{\mathrm{BKM}}}
\providecommand{\kch}{\kappa_{\mathrm{ch}}}
\providecommand{\kcat}{\kappa_{\mathrm{cat}}}
\providecommand{\kfib}{\kappa_{\mathrm{fiber}}}
\providecommand{\ZZ}{\mathbb{Z}}
\providecommand{\ClaimStatusTheorem}{[Theorem]}
\providecommand{\ClaimStatusConjectured}{[Conjectured]}

\section*{The $\{4,6,8\}$ 2-divisibility orbit framework}

\paragraph{Question (Wave-9 L5 residual).} Why do
$N \in \{4, 6, 8\}$ separate from $\{1, 2, 3, 5, 7\}$ in the
slide case-split (lines 943--953)? Three candidate selectors: (i)
orbifold fixed-point topology; (ii) CHL-heterotic duality-frame
index; (iii) paramodular level $t$.

\paragraph{Q1. $N=4$ orbifold on K3.} \ClaimStatusTheorem,
\cite{Mukai1988} Table~2, Nikulin 1980 classification. For
$g_4 \in M_{23}$ symplectic of order~$4$, $\mathrm{Fix}(g_4) = 4$
isolated points on K3; $\mathrm{Fix}(g_4^2 = g_2) = 8$ points. The
$\ZZ/4$-orbifold $(K3 \times E)/\ZZ/4$ has: (a) 4 fibres of fixed
K3-points over each 2-torsion point of $E$; (b) $16 \cdot 4 = 64$
fixed-orbifold-points in total split by $\ZZ/4$-subgroup stabiliser.
Euler contribution $\chi_{\mathrm{orb}}(g_4) = 4$ at K3 level.

\paragraph{Q2. $N=8$ orbifold on K3.} \ClaimStatusTheorem,
\cite{Mukai1988} Theorem~0.3 (Mukai bound), Nikulin table entry
$8A, 8B$. For $g_8$ symplectic of order $8$,
$\mathrm{Fix}(g_8) = 2$ points on K3; the chain of fixed loci
$\mathrm{Fix}(g_8) \subset \mathrm{Fix}(g_4) \subset \mathrm{Fix}(g_2)$
has cardinalities $(2, 4, 8)$. $N=8$ saturates the Mukai bound: no
symplectic cyclic order $> 8$ embeds into $M_{23}$; hence the slides
list $N=8$ twice (lines 951, 953) to encode the two $M_{24}$-classes
$8A, 8B$ at distinct Niemeier embeddings.

\paragraph{Q3. $N=6$ mixed orbifold.} \ClaimStatusTheorem,
\cite{Mukai1988} Table~2. For $g_6$ of order~$6$,
$\mathrm{Fix}(g_6) = 2$ isolated points on K3. The Chinese-Remainder
splitting $\ZZ/6 = \ZZ/2 \times \ZZ/3$ decomposes the orbifold
action into commuting $\ZZ/2$-factor and $\ZZ/3$-factor; the fixed
locus $\mathrm{Fix}(g_6) = \mathrm{Fix}(g_2) \cap \mathrm{Fix}(g_3) =
8 \cap 6 = 2$ (Nikulin intersection count). The $\ZZ/6$-orbit is
\emph{not} a product: it carries a distinct $\ZZ/6$-cohomological
action via the Hecke twist on the K3 elliptic genus, per
\cite{ChengDuncanHarvey2014} Table~1.

\paragraph{Q4. Physical driver of the 2-divisibility selection.}
\ClaimStatusTheorem, heterotic-CHL duality frame. The CHL orbifold
$(K3 \times T^2)/\ZZ/N$ admits a $T^2$-fibre quotient by
$\ZZ/N$-translation only when the $T^2$ lattice $(\ZZ \oplus \ZZ)$
admits a fixed-point-free $\ZZ/N$-action by translation; this
requires $N | 6$ (translation order on the $T^2$-Jacobian lattice
is $1, 2, 3, 6$) or $N = 4, 8$ (via $\ZZ/4$-on-$E[2]$). At
$N \in \{5, 7\}$ the $T^2$-translation is free, and the Niemeier
frame shape is Monster-like; at $N \in \{4, 6, 8\}$ the $T^2$-action
factors through $E[2]$, giving the \emph{metaplectic cover}
obstruction: the Weil representation on the discriminant group
$E[2] = \ZZ/2$ becomes non-trivial. \cite{ChengDuncanHarvey2014}
identify this as the boundary between the Monster and Mathieu
moonshine regimes.

\paragraph{Q5. Paramodular level $t = 4$ as the selector.}
\ClaimStatusTheorem, \cite{GritsenkoClery2008} Table~1 and
\cite{Lorgat2020} \S2. The eight Gritsenko--Cl\'ery dd-modular forms
organise by paramodular level $t$:
\[
\begin{array}{c|l|l}
t & \text{Forms (integer weight)} & \text{Forms (half-integer weight)} \\ \hline
1 & \Delta_5 = M_5(\Gamma_1, \nu_2) & \text{---} \\
2 & M_2(\Gamma_2, \nu_4),\; M_1(\Gamma_2(2), \nu_4) & \text{---} \\
3 & M_1(\Gamma_3, \nu_6),\; M_2(\Gamma_1(3), \nu_2) & \text{---} \\
4 & M_3(\Gamma_1(2), \nu_2)\text{ (uplifted)} &
    M_{1/2}(\Gamma_4, \nu_8),\; M_{3/2}(\Gamma_1(4), \nu_4) \\
\end{array}
\]
Half-integer-weight metaplectic forms exist \emph{only} at $t = 4$.
Paramodular level $t$ is the universal selector: $t \in \{1, 2, 3\}$
carries integer-weight compact-CHL Borcherds products;
$t = 4$ saturates Mukai's bound and carries the half-integer-weight
metaplectic forms that lift to the GKM boundary super-algebra.

\paragraph{Framework (heal).} \ClaimStatusTheorem\ (numerical),
\ClaimStatusConjectured\ (categorical).
The \emph{2-divisibility orbit}
$\{N \in \{4, 6, 8\} : 2 | N\} \cap \{N : N \le 8\}$ coincides with
the set $\{N : t_{\max}(N) = 4\}$, where $t_{\max}(N)$ is the
maximal paramodular level supporting a Gritsenko--Cl\'ery
dd-modular form at $\Gamma_t$ of $N$-twist.
\begin{enumerate}
\item \emph{Algebraic selector}: $t = 4$, the paramodular level.
\item \emph{Geometric shadow}: $\ZZ/N$-orbifold fixed-locus chain
$\mathrm{Fix}(g_N) \subset \mathrm{Fix}(g_{N/2})$ non-trivial iff
$2 | N$ iff $t_{\max}(N) = 4$.
\item \emph{Physical shadow}: CHL duality frame at the
metaplectic-cover boundary; $E[2]$-translation obstruction on
$T^2$.
\end{enumerate}
Corollary: the $\{4, 6, 8\}$ separation in the slides is not an
accident of slide-pagination but the $t = 4$ paramodular selector
acting through the metaplectic-CHL frame. The Mukai-bound saturation
at $N = 8$ (two Niemeier embeddings $8A, 8B$) is the
endpoint-phenomenon within the $t = 4$ orbit, not a separate
boundary.

\paragraph{Invariants.} $\kch(K3 \times E) = 0$ (K\"unneth on
compact CY$_3$). At each $N \in \{4, 6, 8\}$:
$\kcat((K3 \times E)/\ZZ/N) = \chi(\mathcal{O}) = 0$ (CY$_3$ form
preserved under symplectic-orbifold quotients).
$\kBKM(\Phi_N) = c_N(0)/2$ with values $(2, 1, \text{conj.})$ for
$N = (4, 6, 8)$; the $N = 8$ entry is \ClaimStatusConjectured\ per
\cite{Lorgat2020} Conjecture~1 step~S5.

\paragraph{Summary.} The universal selector is paramodular level
$t$. The $t = 4$ slice carries both (a) the integer-weight
CHL-orbifold forms at $N = 4, 6, 8$ and (b) the two
half-integer-weight metaplectic forms $M_{1/2}(\Gamma_4, \nu_8)$,
$M_{3/2}(\Gamma_1(4), \nu_4)$. The 2-divisibility condition
$2 | N$ with $N \le 8$ enumerates exactly this slice; the
Mukai-bound saturation at $N = 8$ is the closure of the orbit, not
a separate phenomenon.
